\section{结论}
在许多格点规范理论应用中需要携带高动量的强子。由于 $n$ 点函数的相对误差随欧氏时间距离指数增长以及基态采样的减弱，高动量以往非常困难
甚至无法实现。在第~\ref{sec:basic} 节
我们引入了一类新的夸克涂抹方法用于构造强子内插算符，
它们应对并显著缓解了这些问题。这些方法的一个具体实现——实现平凡且计算开销极小——是动量 Wuppertal 涂抹，定义见式~\eqref{fermsmear2}。
我们在第~\ref{sec:momtest} 节对其进行了非常成功的检验：仅用 200 次测量便能确定动量分别高达近 $2\units{GeV}$
和 $3\units{GeV}$ 的 $\pi$ 介子与核子能量，见图~\ref{fig:disppi} 与 \ref{fig:dispN}。这两幅图
还包括与传统方法的比较。

在第~\ref{sec:boo} 节我们研究了
引入（附加）洛伦兹助推的可能性~\cite{Roberts:2012tp,DellaMorte:2012xc}。
无论是否包含动量相位因子，
这都比无助推的 Wuppertal 涂抹有一定改进，
可能是因为沿动量方向内插波函数更快的衰减抑制了相位失配。
然而所得结果逊于不加任何 boost 的动量 Wuppertal 涂抹。
我们的结论是：收缩波函数的直觉可能并无根据，
因为强子在真实时间里运动而不在虚（欧氏）
时间里运动。

迭代方法
（无论是否动量涂抹）随着连续极限 $a\rightarrow 0$
的逼近而苦于高迭代次数
$n\propto a^{-2}$，此外还有格点数目增加带来的朴素体积因子。因此，在第~\ref{sec:improve} 节
我们引入了其他非迭代（动量）涂抹方法，
它们可能更适合小的格距，
并且也允许构造非高斯形状。
我们近期将使用这些方法。

实现远大于相关强子质量的动量在若干现代应用中具有根本重要性，
例如直接确定（准）分布
振幅~\cite{Aglietti:1998ur,Abada:2001if,Braun:2007wv}、
（准）（广义）部分子分布~\cite{Ji:2013dva,Ji:2015jwa,Lin:2014zya,Alexandrou:2015rja}
以及横动量分布的矩~\cite{Hagler:2009mb,Musch:2010ka,Musch:2011er,Engelhardt:2015xja}。
新方法使我们能够（以非常有限的计算资源）提取高达动量 $\vect{p}^2\approx 7.9\units{GeV}^2$ 的核子质量。
这清楚地使上述可观测量适合现实环境下的格点模拟。对三点函数而言，这些方法
甚至可能允许虚度
$Q^2=4\vect{p}^2\approx 30\units{GeV}^2$：把源动量
$-\vect{p}$ 切换为汇动量 $+\vect{p}$。
用新涂抹计算这些量的工作视可观测量而定正在进行或已在计划之中。\\
